% Text overlay only; model artwork is retained without modification.
\begin{tikzpicture}[x=1mm,y=1mm,text=BookInk]
\path[use as bounding box] (0,0) rectangle (157.4,71.50698924731182);
\node[inner sep=0pt,outer sep=0pt,align=left,anchor=west,font=\figurefont\fontsize{8}{10}\selectfont,text width=50.368mm] at (6.2960,66.8590) {A  递推记忆};
\node[inner sep=0pt,outer sep=0pt,align=left,anchor=west,font=\figurefont\fontsize{8}{10}\selectfont,text width=53.516000000000005mm] at (83.4220,66.8590) {B  历史注意力};
\node[inner sep=0pt,outer sep=0pt,align=center,anchor=center,font=\figurefont\fontsize{8}{10}\selectfont,text width=25.184mm] at (16.0548,61.1385) {\mbox{旧记忆$\bm z_t$}};
\node[inner sep=0pt,outer sep=0pt,align=center,anchor=center,font=\figurefont\fontsize{8}{10}\selectfont,text width=15.740000000000002mm] at (37.7760,44.1198) {\mbox{递推$f_{\bm\theta}$}};
\node[inner sep=0pt,outer sep=0pt,align=center,anchor=center,font=\figurefont\fontsize{8}{10}\selectfont,text width=25.184mm] at (59.6546,61.1385) {\mbox{新记忆$\bm z_{t+1}$}};
\node[inner sep=0pt,outer sep=0pt,align=center,anchor=east,font=\figurefont\fontsize{8}{10}\selectfont,text width=28.332mm] at (28.3320,26.7436) {动作与反馈$(a_t,r_t,o_{t+1})$};
\node[inner sep=0pt,outer sep=0pt,align=center,anchor=center,font=\figurefont\fontsize{8}{10}\selectfont,text width=14.166mm] at (59.6546,22.0242) {策略};
\node[inner sep=0pt,outer sep=0pt,align=center,anchor=center,font=\figurefont\fontsize{8}{10}\selectfont,text width=34.628mm] at (94.4400,57.2056) {追加动作与反馈};
\node[inner sep=0pt,outer sep=0pt,align=right,anchor=east,font=\figurefont\fontsize{8}{10}\selectfont,text width=25.184mm] at (77.1260,34.7524) {\mbox{历史$I_{t+1}$}};
\node[inner sep=0pt,outer sep=0pt,align=center,anchor=center,font=\figurefont\fontsize{8}{10}\selectfont,text width=12.592mm] at (126.5496,29.1749) {注意力};
\node[inner sep=0pt,outer sep=0pt,align=center,anchor=center,font=\figurefont\fontsize{8}{10}\selectfont,text width=23.61mm] at (145.1228,61.1385) {\mbox{编码$\bm z_{t+1}$}};
\node[inner sep=0pt,outer sep=0pt,align=center,anchor=center,font=\figurefont\fontsize{8}{10}\selectfont,text width=14.166mm] at (145.2802,22.0242) {策略};
\node[inner sep=0pt,outer sep=0pt,align=center,anchor=center,font=\FigureNoteFont\fontsize{7.5}{9.5}\selectfont,text width=42.498000000000005mm] at (26.7580,10.9406) {预热前缀：重建记忆};
\node[inner sep=0pt,outer sep=0pt,align=center,anchor=center,font=\FigureNoteFont\fontsize{7.5}{9.5}\selectfont,text width=78.7mm] at (102.3100,10.9406) {训练后缀：计算损失};
\node[inner sep=0pt,outer sep=0pt,align=center,anchor=center,font=\FigureNoteFont\fontsize{7.5}{9.5}\selectfont,text width=147.956mm] at (78.7000,2.6458) {真实回合起点：重置记忆；片段损失起点：保持已重建的记忆};
\end{tikzpicture}
